\clearpage
\section{Problem 2: Concurrency, Availability, dan Graceful Degradation}

\subsection{Rancangan Solusi dan Mekanisme}
Permasalahan kedua (P2) berfokus pada pencegahan tiga titik kegagalan utama dalam sistem terdistribusi saat berhadapan dengan beban tinggi dan gangguan penyedia data eksternal:
\begin{enumerate}
    \item \textbf{Pencegahan Head-of-Line Blocking:} Aggregator mengisolasi pemanggilan BMKG dan PVMBG ke \textit{goroutine} asinkron mandiri (\textit{timeout} 2 detik BMKG, 5 detik PVMBG). Kelambatan ekstrem atau penundaan (3 detik) PVMBG tidak menghambat penarikan data seismik BMKG.
    \item \textbf{Pencegahan Resource Exhaustion:} Pada lapisan HTTP publik (\path{client-api}), konkurensi dibatasi oleh penyangga HTTP Go dan koneksi basis data \path{pgxpool} (\texttt{MaxConns = 50}), mencegah kehabisan soket dan memori saat lonjakan beban.
    \item \textbf{Penerapan Graceful Degradation:} Aggregator menerapkan pelacak keandalan berbasis \textit{Circuit Breaker}. Jika PVMBG gagal beruntun ($\ge 3$) atau menghasilkan HTTP 503, sistem mengaktifkan status degradasi (\texttt{is\_degraded = true}). Respons klien tetap HTTP 200 berisi data BMKG terbaru dipadu \textit{snapshot} historis PVMBG beserta metadata penanda status.
\end{enumerate}

\subsection{Alasan Keputusan dan Trade-Off}
Dalam menangani konkurensi dan degradasi, kelompok mengevaluasi trade-off berikut:
\begin{itemize}
    \item \textbf{Pola Fail-Fast (HTTP 503 Langsung ke Klien):} Menolak permintaan klien secara instan saat salah satu sumber mati. Ditolak karena menghilangkan kegunaan data seismik bagi tim lapangan saat hanya sumber vulkanik yang terganggu.
    \item \textbf{Stale-Data Fallback dengan Penanda Kebaruan (Solusi Terpilih):} Sistem mengutamakan ketersediaan (\textit{Availability} per teorema CAP). Klien menerima data yang ada disertai transparansi status melalui field \path{meta.degraded = true} dan stempel waktu \path{meta.stale_since}.
\end{itemize}

\clearpage
\subsection{Batasan yang Disadari}
\begin{itemize}
    \item Pada kondisi degradasi berkepanjangan, data vulkanik yang disajikan semakin usang (\textit{stale}). Sistem Milestone 1 mengandalkan pemahaman klien terhadap atribut \path{meta.stale_since} dan belum mengimplementasikan estimasi interpolasi bencana sintetis.
    \item Pembatasan laju permintaan (\textit{rate limiting}) pada Milestone 1 ditangani pada level kapasitas koneksi server (\textit{connection backlog}), belum menggunakan token bucket terdistribusi Redis.
\end{itemize}

\subsection{Bukti Penyelesaian}
Verifikasi empiris dilakukan melalui dua skrip utama:
\begin{enumerate}
    \item \textbf{Uji Isolasi Latensi \& Simulasi Outage (\texttt{make demo-p2}):}
    PVMBG dimatikan melalui \path{POST /admin/outage} (\texttt{is\_outage = true}). Klien API tetap melayani permintaan dengan status HTTP 200 dan menghasilkan payload terdegradasi:

\noindent\begin{minipage}{\linewidth}
\begin{lstlisting}[language=json, caption={Respons Klien API dalam Keadaan Terdegradasi (meta.degraded = true)}]
{
  "data": [ ... ],
  "meta": {
    "total": 1562, "limit": 5, "offset": 0,
    "degraded": true,
    "stale_since": "2026-10-08T09:45:21.5830728Z"
  }
}
\end{lstlisting}
\end{minipage}
    Sementara itu, pemanggilan langsung ke \path{mock-bmkg} (:8081) tetap merespons dengan HTTP 200 dan latensi p95 jauh di bawah ambang 300 ms (\path{docs/evidence/p2/demo_p2.log}).

    \item \textbf{Uji Beban Konkurensi Berkelanjutan (\texttt{make load-test}):}
    Pengujian beban menggunakan k6 dijalankan dengan parameter 50 Virtual Users (VUs) secara simultan selama 60 detik penuh terhadap titik akhir \texttt{GET http://localhost:8080/api/v1/hazards} (\path{docs/evidence/p2/load_test.log}):
    \begin{itemize}
        \item \textbf{Total Permintaan Terlayani:} 29.514 permintaan ($\approx 491$ req/s).
        \item \textbf{Ambang Batas Latensi \& Error Rate:} Latensi p95: 1,75 ms ($< 300$ ms); Latensi p99: 3,07 ms ($< 500$ ms); Error Rate: 0,43\% ($< 1,0\%$) --- seluruh kriteria evaluasi beban terbukti \mbox{\textbf{MEMENUHI}}.
    \end{itemize}
\end{enumerate}

\clearpage
\subsection{Jawaban Butir Spesifikasi (Di Laporan)}
\textbf{Pertanyaan:} \textit{Jelaskan model konkurensi yang dipilih dan bagaimana model tersebut mencegah ketiga kegagalan di atas. Sertakan konfigurasi load testing beserta tooling yang dipakai.}

\textbf{Jawaban:}
\begin{enumerate}
    \item \textbf{Model Konkurensi Terpilih:} Sistem mengadopsi model konkurensi berbasis \textit{Communicating Sequential Processes (CSP)} yang diwujudkan melalui \textit{goroutines} dan \textit{buffered channels} di Go:
    \begin{itemize}
        \item \textit{Head-of-line blocking} dicegah dengan memisahkan goroutine penarik BMKG dan PVMBG secara asinkron; penundaan pada satu socket TCP upstream tidak menghambat pemrosesan socket lainnya.
        \item \textit{Resource exhaustion} dicegah melalui pembatasan ukuran \textit{connection pool} PostgreSQL (\texttt{pgxpool.Config.MaxConns = 50}) dan pemanfaatan konkurensi ringan goroutine (hanya memakan memori $\approx 2$ KB per goroutine dibanding utas OS standar 2--8 MB).
        \item \textit{Ketiadaan degradasi terkendali} dicegah dengan pencatatan state kesehatan penyedia data; saat kegagalan terdeteksi, Aggregator beralih menyajikan snapshot terakhir basis data kanonis dengan flag \texttt{meta.degraded = true}.
    \end{itemize}
    \item \textbf{Konfigurasi Load Testing:}
    \begin{itemize}
        \item \textbf{Tooling:} Grafana k6 v0.50+ dieksekusi secara otomatis melalui Docker atau biner lokal (\path{scripts/load_test.js}).
        \item \textbf{Beban \& Durasi:} 50 Virtual Users (VUs) selama 60 detik konstan (\textit{sustained load}).
        \item \textbf{Autentikasi:} Setiap VU melakukan autentikasi otomatis ke Auth Service dan menyertakan header \texttt{Authorization: Bearer <token>} pada setiap iterasi.
    \end{itemize}
\end{enumerate}
